\documentclass[11pt,a4paper]{article}
\usepackage[utf8]{inputenc}
\usepackage[T1]{fontenc}
\usepackage{lmodern}
\usepackage[italian]{babel}
\usepackage{amsmath, amssymb, mathtools}
\usepackage{graphicx}
\usepackage[margin=2.5cm]{geometry}
\usepackage{xcolor}
\usepackage{hyperref}
\usepackage{algorithm}
\usepackage[noend]{algpseudocode} % noend per non mostrare "end if", "end for", ecc.
\usepackage{booktabs} % Per linee di tabella più professionali
\usepackage{multirow} % Per celle che si estendono su più righe
\usepackage{siunitx}  % Per allineare i numeri per punto decimale e formattazione
\usepackage{enumitem} % Per personalizzare le liste
\usepackage{float} % <--- AGGIUNGI QUESTO PACCHETTO

\usepackage{listings}
\usepackage{xcolor}

% Configura l'aspetto dei blocchi di codice
\lstset{
    basicstyle=\ttfamily\small,
    commentstyle=\color{green!50!black},
    keywordstyle=\color{blue},
    stringstyle=\color{red},
    breaklines=true,
    breakatwhitespace=true,
    tabsize=4,
    showstringspaces=false,
    frame=single
}

\hypersetup{
    colorlinks=true,
    linkcolor=blue,
    filecolor=magenta,      
    urlcolor=cyan,
    pdftitle={Spiegazione di TIMO e TIMO-S},
    pdfpagemode=FullScreen,
}


\title{Text and Image Are Mutually Beneficial (TIMO) per il Few-Shot Classification Training-Free con CLIP e dataset usato}
\author{Basato sul paper di Yayuan Li, Jintao Guo, Lei Qi, Wenbin Li, Yinghuan Shi}
\date{\today}

\newcommand{\Ft}{\mathbf{F}_\text{t}}      % Text features
\newcommand{\Fv}{\mathbf{F}_\text{v}}      % Visual features (support)
\newcommand{\Wt}{\mathbf{W}_\text{t}}      % Text prototypes
\newcommand{\Wv}{\mathbf{W}_\text{v}}      % Visual prototypes (support)
\newcommand{\fv}{\mathbf{f}_\text{v}}      % Query visual feature
\newcommand{\Figt}{\mathbf{F}_\text{IGT}}  % IGT features
\newcommand{\Ftgi}{\mathbf{F}_\text{TGI}}  % TGI features
\newcommand{\R}{\mathbb{R}}

\begin{document}
\maketitle
\begin{abstract}
Questo documento descrive TIMO (Text-Image Mutual guidance Optimization) e la sua variante migliorata TIMO-S, metodi proposti per migliorare il Few-Shot Learning (FSL) training-free utilizzando il modello CLIP (Contrastive Language-Image Pretraining). L'obiettivo è affrontare due problemi chiave derivanti dalla modellazione indipendente delle modalità testuale e visiva: (1) un severo disallineamento anomalo nella modalità immagine e (2) una qualità variabile dei prompt testuali generati. TIMO introduce un meccanismo di guida reciproca tra testo e immagine per mitigare questi problemi.
\end{abstract}

\newpage
\tableofcontents
\newpage

\section{Introduzione: Problemi nel Few-Shot Learning con CLIP}
CLIP ha dimostrato notevoli capacità di generalizzazione zero-shot e si è affermato come una base potente per il Few-Shot Learning (FSL). Tuttavia, molti metodi FSL basati su CLIP che non richiedono addestramento aggiuntivo (training-free) tendono a trattare le modalità testuale e visiva in modo indipendente. Questo approccio presenta due limiti principali:

\begin{enumerate}
    \item \textbf{Disallineamento Anomalo nella Modalità Immagine (Anomalous Match):}
    Durante il pre-addestramento di CLIP, possono emergere delle "scorciatoie" o bias che portano a una somiglianza elevata tra rappresentazioni di immagini appartenenti a classi diverse. Questo problema è particolarmente pronunciato in FSL, dove il numero limitato di campioni di supporto può portare il modello a fare previsioni errate con alta confidenza. La modalità visiva, essendo più ricca di informazioni rispetto a quella testuale, è più suscettibile a queste imperfezioni intrinseche di CLIP.

    \item \textbf{Qualità Variabile dei Prompt Testuali Generati:}
    I metodi training-free spesso si affidano a prompt testuali generati automaticamente (ad esempio, da Large Language Models) per descrivere le classi. Tuttavia, la qualità di questi prompt può variare significativamente. Alcuni prompt possono essere poco descrittivi, semanticamente distanti dalle caratteristiche visive delle immagini, o addirittura fuorvianti. Questo divario semantico tra le rappresentazioni testuali e visive può portare a classificatori sub-ottimali, specialmente quando si considerano più prompt per classe. Il paper evidenzia come i prompt peggiori possano degradare significativamente le prestazioni.
\end{enumerate}

Per affrontare congiuntamente questi problemi, TIMO propone un meccanismo di \textbf{guida reciproca (mutual guidance)} tra testo e immagine.

\section{Preliminari: CLIP e Classificazione Few-Shot}
CLIP impara uno spazio di embedding congiunto in cui le rappresentazioni di immagini e testi semanticamente simili sono vicine.
\begin{itemize}
    \item Data una query immagine, CLIP ne estrae la feature visiva $\fv \in \R^D$.
    \item Per ogni classe $i$ (su $N$ classi), si generano $P$ prompt testuali (es. "una foto di una [CLASSE]"). Questi vengono codificati da CLIP per ottenere le feature testuali $\Ft \in \R^{N \times P \times D}$. Solitamente, si calcola una media dei $P$ prompt per ottenere un prototipo testuale per classe $\Wt \in \R^{N \times D}$.
    \item In FSL, si dispone di $K$ immagini di supporto per classe. Queste vengono codificate per ottenere le feature visive di supporto $\Fv \in \R^{N \times K \times D}$, da cui si possono derivare prototipi visivi per classe $\Wv \in \R^{N \times D}$.
\end{itemize}
La classificazione zero-shot standard in CLIP calcola i logits come la similarità coseno tra $\fv$ e $\Wt$.
I metodi FSL training-free spesso combinano i logits testuali (zero-shot) con i logits derivati da un classificatore basato sulle immagini di supporto (es. $\text{Classifier}_{\text{image}}(\fv)$), tipicamente tramite una somma pesata:
\[ \text{logits} = \text{logits}_{\text{text}} + \alpha \cdot \text{logits}_{\text{image}} \]
Tuttavia, come discusso, la costruzione indipendente di questi classificatori è sub-ottimale.

\section{TIMO: Ottimizzazione a Guida Mutua Testo-Immagine}
TIMO introduce due componenti chiave: Image-Guided Text (IGT) e Text-Guided Image (TGI), che lavorano sinergicamente.

\subsection{Text-Guided Image (TGI): Guida Testuale per l'Immagine}
\textbf{Obiettivo:} Mitigare il problema del disallineamento anomalo nella modalità immagine, sfruttando le rappresentazioni testuali, generalmente più robuste.
\textbf{Meccanismo:}
\begin{enumerate}
    \item Per ogni classe $i$ e per ciascuno dei suoi $P$ prompt testuali (feature $F_t^{i,p}$), si calcola una matrice di pesi $s^{i,p}$ basata sulla similarità coseno tra il prompt $F_t^{i,p}$ e il prototipo visivo della classe $W_v^i$:
    \[ s^{i,p} = \frac{F_t^{i,p} \cdot (W_v^i)^T}{\|F_t^{i,p}\| \|W_v^i\|} \]
    Questo peso indica quanto un prompt testuale è rilevante per le caratteristiche visive della classe.
    \item Si introduce un iperparametro $\beta$ che controlla il numero di prompt testuali da considerare per la guida. I pesi $s^{i,p}$ vengono utilizzati per selezionare/pesare i prompt più rilevanti. Sia $S \in \R^{N \times P}$ la matrice di questi pesi (eventualmente con zero per i prompt scartati).
    \item Le feature di supporto visive $\Fv$ vengono combinate con le feature testuali pesate $F_t \odot S$ (prodotto element-wise, o una selezione basata su $\beta$) per ottenere una rappresentazione arricchita $\Ftgi$:
    \[ \Ftgi = \text{Concat}(\Fv, \Ft \odot S) \in \R^{N \times (K+\beta') \times D} \]
    dove $\beta'$ è il numero effettivo di prompt testuali usati dopo la selezione.
    \item Un classificatore per immagini, $\text{Classifier}_{\text{TGI}}$, viene costruito utilizzando queste feature $\Ftgi$ (ad esempio, seguendo la logica di Tip-Adapter o GDA-CLIP, che il paper chiama `ImageClassifierBuilder`).
    \item I logits TGI per l'immagine di query $\fv$ sono: $\text{logits}_{\text{TGI}} = \text{Classifier}_{\text{TGI}}(\fv)$.
\end{enumerate}
In sostanza, TGI "raffina" le feature visive di supporto includendo le informazioni testuali più allineate, con l'obiettivo di creare un classificatore di immagini più robusto ai disallineamenti.

\subsection{Image-Guided Text (IGT): Guida Visuale per il Testo}
\textbf{Obiettivo:} Rettificare la qualità variabile dei prompt testuali, utilizzando le rappresentazioni delle immagini come guida per selezionare o pesare i prompt in modo più efficace.
\textbf{Meccanismo:}
Dal punto di vista dell'ensemble learning, un buon classificatore IGT deve soddisfare due obiettivi:
\begin{enumerate}
    \item \textbf{Target I:} Classificare correttamente i dati.
    \item \textbf{Target II:} Mantenere l'individualità (diversità) delle informazioni originali dei prompt.
\end{enumerate}
Per fare ciò:
\begin{enumerate}
    \item Per ogni classe $i$, si cerca un vettore di pesi $r^i \in \R^P$ per i $P$ prompt testuali $F_t^i$ che massimizzi l'allineamento con il prototipo visivo $W_v^i$ della classe (dopo normalizzazione L2 di tutte le feature):
    \[ \max_{r^i} (r^i)^T F_t^i (W_v^i)^T \]
    soggetto al vincolo che la norma di $r^i$ sia controllata da un iperparametro $\gamma$: $\|r^i\| = \gamma$. Questo vincolo aiuta a preservare la diversità dei prompt (Target II).
    \item La soluzione ottimale per $r^i$ è proporzionale a $F_t^i (W_v^i)^T$. Questi pesi $r^i$ vengono poi normalizzati (ad es. con Softmax) per ottenere una distribuzione $R^i$ sui prompt. Sia $R \in \R^{N \times P}$ la matrice di questi pesi combinati per tutte le classi.
    \item Le feature testuali originali $\Ft$ vengono trasformate utilizzando questi pesi per ottenere le feature testuali rettificate $\Figt$:
    \[ \Figt = \Ft R^T \in \R^{N \times D} \quad (\text{o } \Figt = \text{einsum("np, npd -> nd", R, Ft)}) \]
    \item Un classificatore testuale, $\text{Classifier}_{\text{IGT}}$, viene costruito utilizzando queste feature $\Figt$ (che il paper chiama `TextClassifierBuilder`).
    \item I logits IGT per l'immagine di query $\fv$ sono: $\text{logits}_{\text{IGT}} = \text{Classifier}_{\text{IGT}}(\fv)$.
\end{enumerate}
IGT quindi ri-pesa o combina linearmente i prompt testuali per ogni classe in modo che siano più semanticamente allineati con le caratteristiche visive di quella classe.

\subsection{Combinazione in TIMO}
TIMO combina i logits ottenuti dai classificatori TGI e IGT tramite una somma pesata, controllata da un iperparametro $\alpha$:
\[ \text{logits}_{\text{TIMO}} = \text{logits}_{\text{TGI}} + \alpha \cdot \text{logits}_{\text{IGT}} \]
L'iperparametro $\alpha$ bilancia il contributo delle due modalità guidate. Viene determinato tramite grid search su un set di validazione.

\section{TIMO-S: Versione Migliorata di TIMO}
TIMO-S è una variante potenziata di TIMO. Mentre TIMO ottimizza principalmente $\alpha$, TIMO-S estende l'ottimizzazione includendo una grid search anche per gli iperparametri $\beta$ (per TGI, che controlla il grado di guida testuale per l'immagine) e $\gamma$ (per IGT, che controlla il grado di guida visiva per il testo) su un set di validazione.
Questo permette di trovare una configurazione più robusta e ottimale per il grado di guida reciproca, portando a prestazioni ulteriormente migliorate. Il paper riporta che TIMO-S supera i metodi training-required esistenti con un costo computazionale significativamente inferiore.

\section{Dataset Utilizzato}
\label{sec:dataset}
Per la valutazione di TIMO e TIMO-S, è stato impiegato un dataset derivato da una collezione più ampia di opere d'arte, comunemente nota come ArtGraph. La creazione del dataset specifico per gli esperimenti ha seguito una logica precisa per garantire la pertinenza e la gestibilità nell'ambito del few-shot learning.

\subsection{Composizione del Dataset Originale}
Il dataset di partenza è costituito da un vasto insieme di immagini di opere d'arte, ciascuna associata a metadati quali il nome dell'artista e lo stile artistico predominante. Queste informazioni sono tipicamente raccolte in un formato tabellare (un file CSV), dove ogni riga corrisponde a un'opera d'arte.

\subsection{Creazione del Nuovo Dataset per la Sperimentazione}
La costruzione del dataset utilizzato per gli esperimenti di classificazione few-shot, con l'obiettivo di classificare le opere per artista, ha comportato diversi passaggi di selezione e strutturazione:
\begin{enumerate}
    \item \textbf{Selezione Basata sullo Stile e Numero di Opere:}
    Inizialmente, sono stati identificati un numero limitato dei stili artistici più rappresentati nel dataset originale (i primi 10 stili per numero di opere).
    \item \textbf{Selezione degli Artisti per Stile:}
    All'interno di ciascuno di questi stili principali, sono stati selezionati un certo numero di artisti. Un criterio fondamentale per la selezione di un artista è stato il numero di opere possedute \textit{all'interno di quello specifico stile}, che doveva rientrare in un intervallo predefinito (tra 30 e 50 opere). Questo assicura che ci siano abbastanza campioni per artista per creare split significativi, ma non così tanti da sbilanciare eccessivamente il dataset o allontanarsi troppo dallo scenario few-shot.
    \item \textbf{Esclusione degli Altri Artisti:}
    È importante sottolineare che solo gli artisti che soddisfacevano i criteri di cui sopra (appartenenza a uno stile top e numero di opere adeguato in quello stile) sono stati inclusi nel dataset finale. Tutti gli altri artisti e le loro opere sono stati esclusi dalla sperimentazione.
    \item \textbf{Divisione in Set di Support, Validazione e Query:}
    Per ciascun artista selezionato, le sue opere (appartenenti allo stile per cui è stato scelto) sono state suddivise percentualmente in tre insiemi distinti:
    \begin{itemize}
        \item \textbf{Support set:} Una porzione maggioritaria delle opere (80\%), utilizzata per "mostrare" al modello le caratteristiche dell'artista durante la fase di adattamento (support set nel few-shot learning).
        \item \textbf{Validation set:} Una piccola porzione (10\%), usata per l'ottimizzazione degli iperparametri del modello (come $\alpha$, $\beta$, $\gamma$ in TIMO/TIMO-S).
        \item \textbf{Query set:} La restante porzione (10\%), utilizzata per valutare le prestazioni finali del modello su dati mai visti (query set nel few-shot learning).
    \end{itemize}
    \item \textbf{Strutturazione Finale:}
    Il dataset risultante è stato organizzato in una struttura di directory standard, con le immagini raggruppate per artista e file di etichette separati che definiscono gli split di training, validazione e test. A ogni artista selezionato è stato assegnato un ID numerico univoco.
\end{enumerate}
Questa metodologia di creazione del dataset mira a costruire un benchmark focalizzato e bilanciato per la classificazione few-shot di artisti, basandosi su una selezione ragionata degli stili e del numero di opere per artista.

\subsection{Ruolo dei Prompt Testuali Generati da LLM}
Nel contesto di modelli multimodali come CLIP, e quindi anche per TIMO, i \textbf{prompt testuali} giocano un ruolo cruciale. Termini come \texttt{Prompt\_CuPL}, \texttt{Prompt\_Waffle}, \texttt{Prompt\_DCLIP}, ecc., si riferiscono a insiemi di descrizioni testuali generate per ogni classe del dataset (in questo caso, per ogni artista).
\begin{itemize}
    \item \textbf{Generazione tramite Grandi Modelli Linguistici (LLM):} Questi prompt sono tipicamente prodotti utilizzando LLM. Per ogni artista, l'LLM viene istruito a generare molteplici frasi descrittive. Ad esempio, per un artista specifico, i prompt potrebbero includere varianti come "un dipinto di [Nome Artista]", "opere d'arte nello stile distintivo di [Nome Artista]", "arte creata da [Nome Artista]", ecc.
    \item \textbf{Scopo dei Prompt:}
    \begin{enumerate}
        \item \textbf{Creare Rappresentazioni Testuali Ricche:} Fornire diverse descrizioni testuali per ogni artista aiuta a catturare un'ampia gamma di sfumature semantiche associate al suo lavoro e stile.
        \item \textbf{Migliorare la Generalizzazione di CLIP:} CLIP apprende ad associare immagini e testi. Utilizzando un insieme variegato di prompt per rappresentare testualmente un artista, si può ottenere una rappresentazione vettoriale (un "prototipo testuale") più robusta e generalizzabile per quell'artista.
        \item \textbf{Supporto al Few-Shot/Zero-Shot Learning:} In scenari con pochi (o nessun) esempi di addestramento per artista, avere prototipi testuali di alta qualità è fondamentale. I prompt generati da LLM permettono di definire le classi al modello anche senza un gran numero di immagini di esempio.
        \item \textbf{Guida per l'Ottimizzazione (come in TIMO):} Metodi come TIMO utilizzano questi prompt come punto di partenza. La componente IGT (Image-Guided Text) di TIMO, ad esempio, cerca di raffinare o pesare questi prompt testuali utilizzando le informazioni provenienti dalle immagini di supporto, al fine di migliorare l'allineamento tra le modalità visiva e testuale e, di conseguenza, le prestazioni di classificazione.
    \end{enumerate}
\end{itemize}
In sintesi, i prompt generati da LLM forniscono la necessaria controparte testuale che permette a modelli come CLIP di comprendere e classificare le immagini degli artisti, e costituiscono una componente essenziale per le strategie di adattamento e ottimizzazione impiegate in metodi avanzati come TIMO.

\section{Risultati}

\subsection{Tabelle dei Risultati}
Vengono presentate di seguito le tabelle riassuntive delle performance medie ottenute sul dataset ArtGraph, valutando Accuratezza, Precisione Macro, Recall Macro e F1-Score Macro.

\begin{table}[H]
    \centering
    \caption{Risultati Medi di Accuratezza (\%) su ArtGraph. Ogni valore rappresenta la media delle accuratezze ottenute su 10 diversi seed per la specifica combinazione di Backbone, Numero di Shots e Modello di Adattamento.}   
    \sisetup{
      round-mode=places,
      round-precision=2,
      output-decimal-marker={,},
      table-format=2.2 
    }
    \resizebox{\textwidth}{!}{%
    \begin{tabular}{@{}lcSSSSS@{}}
    \toprule
    Backbone & Shots & {\rotatebox{90}{APE}} & {\rotatebox{90}{GDA\_CLIP}} & {\rotatebox{90}{TIMO}} & {\rotatebox{90}{TIMO\_S}} & {\rotatebox{90}{Tip\_Adapter}} \\
    \midrule
    \multirow{5}{*}{RN101}
    & 1  & 42.354740061162076 & 41.681957186544345 & 45.90214067278288 & 45.90214067278288 & 38.53211009174312 \\
    & 2  & 47.73700305810398 & 46.88073394495413 & 47.522935779816514 & 49.20489296636086 & 41.13149847094802 \\
    & 4  & 52.41590214067278 & 55.657492354740064 & 57.88990825688073 & 58.25688073394495 & 45.50458715596331 \\
    & 8  & 57.30886850152905 & 64.06727828746178 & 64.61773700305811 & 65.01529051987768 & 49.78593272171254 \\
    & 16 & 59.1131498470948 & 71.28440366972477 & 71.131498470948 & 71.31498470948011 & 54.189602446483185 \\
    \midrule
    \multirow{5}{*}{RN50}
    & 1  & 44.31192660550458 & 41.406727828746185 & 43.883792048929664 & 45.963302752293586 & 38.654434250764524 \\
    & 2  & 47.09480122324159 & 47.217125382263 & 47.88990825688073 & 48.65443425076453 & 40.73394495412844 \\
    & 4  & 53.883792048929664 & 55.535168195718654 & 56.48318042813455 & 57.2782874617737 & 44.95412844036697 \\
    & 8  & 58.103975535168196 & 63.608562691131496 & 65.25993883792049 & 65.07645259938838 & 48.440366972477065 \\
    & 16 & 62.171253822629964 & 73.05810397553516 & 73.36391437308869 & 73.45565749235473 & 52.90519877675841 \\
    \midrule
    \multirow{5}{*}{ViT-B/16}
    & 1  & 50.8868501529052 & 48.92966360856269 & 52.38532110091743 & 53.57798165137615 & 47.06422018348624 \\
    & 2  & 54.77064220183486 & 55.535168195718654 & 54.434250764526 & 56.20795107033639 & 49.14373088685015 \\
    & 4  & 60.21406727828746 & 63.149847094801224 & 64.25076452599389 & 64.74006116207951 & 52.59938837920489 \\
    & 8  & 63.82262996941896 & 70.73394495412843 & 71.40672782874618 & 71.34556574923548 & 58.1651376146789 \\
    & 16 & 68.47094801223241 & 77.70642201834862 & 78.10397553516819 & 78.01223241590215 & 62.69113149847095 \\
    \midrule
    \multirow{5}{*}{ViT-B/32}
    & 1  & 45.50458715596331 & 43.36391437308869 & 47.522935779816514 & 48.13455657492355 & 41.37614678899082 \\
    & 2  & 48.13455657492355 & 47.3394495412844 & 49.48012232415902 & 49.357798165137616 & 42.50764525993884 \\
    & 4  & 54.98470948012233 & 56.666666666666664 & 59.32721712538226 & 59.44954128440368 & 45.90214067278288 \\
    & 8  & 59.87767584097859 & 63.45565749235474 & 64.46483180428135 & 64.434250764526 & 50.0 \\
    & 16 & 62.8440366972477 & 70.42813455657492 & 71.28440366972477 & 71.28440366972477 & 56.91131498470948 \\
    \bottomrule
    \end{tabular}%
    } % Fine resizebox
    \label{tab:artgraph_results_avg_accuracy}
\end{table}

\begin{table}[H]
    \centering
    \caption{Risultati Medi di Precisione Macro (\%) su ArtGraph. Ogni valore rappresenta la media dei punteggi di Precisione Macro ottenuti su 10 diversi seed.}
    \sisetup{
      round-mode=places,
      round-precision=2,
      output-decimal-marker={,},
      table-format=2.2 
    }
    \resizebox{\textwidth}{!}{%
    \begin{tabular}{@{}lcSSSSS@{}}
    \toprule
    Backbone & Shots & {\rotatebox{90}{APE}} & {\rotatebox{90}{GDA\_CLIP}} & {\rotatebox{90}{TIMO}} & {\rotatebox{90}{TIMO\_S}} & {\rotatebox{90}{Tip\_Adapter}} \\
    \midrule
    \multirow{5}{*}{RN101}
    & 1  & 43.05135379900868 & 40.255692117024786 & 46.962843869093874 & 45.897433485421125 & 35.89954348624519 \\
    & 2  & 50.0695614656528 & 48.5759032360393 & 49.0394007034632 & 50.435037933199695 & 38.65352921352026 \\
    & 4  & 56.03590620819239 & 59.27504234438735 & 61.51508474164724 & 61.757158350908355 & 44.32104039688994 \\
    & 8  & 60.83930218236285 & 66.49949616355866 & 67.57454004329006 & 67.76104797979798 & 49.85838001468845 \\
    & 16 & 62.74830031080031 & 73.71490575396825 & 73.55195932539682 & 73.5905257936508 & 55.84780397674651 \\
    \midrule
    \multirow{5}{*}{RN50}
    & 1  & 44.419927192905135 & 38.86063532074466 & 44.3172037147221 & 44.472965965564654 & 30.849328670644944 \\
    & 2  & 47.769854073622454 & 47.89951399873275 & 48.42613078657406 & 48.7832323756037 & 32.99438541979177 \\
    & 4  & 54.925460304838204 & 58.23556564962816 & 59.137366973304474 & 59.7717511076886 & 39.15567157640977 \\
    & 8  & 58.93553408397158 & 67.18288297010723 & 68.42943948412699 & 68.2510333994709 & 45.25988407876953 \\
    & 16 & 63.6503351972102 & 76.11747685185185 & 76.58767823611574 & 76.51931980056979 & 51.89539511414512 \\
    \midrule
    \multirow{5}{*}{ViT-B/16}
    & 1  & 53.06164146830265 & 48.31639212081022 & 53.38168602231102 & 54.10456094831095 & 43.101542041424494 \\
    & 2  & 57.12861586495004 & 57.81490153365153 & 55.799199642949645 & 57.718036593036594 & 47.71252612286378 \\
    & 4  & 62.98829787111036 & 66.03974237567988 & 66.32724942881194 & 67.12920499639249 & 52.45289975565227 \\
    & 8  & 66.83949990981242 & 73.74673445767196 & 74.20258763227513 & 74.08316798941799 & 59.642196440631515 \\
    & 16 & 71.46413915945166 & 80.36995701058201 & 80.62251984126983 & 80.42427248677248 & 65.81936235061235 \\
    \midrule
    \multirow{5}{*}{ViT-B/32}
    & 1  & 44.49598720485669 & 40.33418190265617 & 46.15219719516595 & 46.50982929107929 & 37.24655006111807 \\
    & 2  & 49.20309281200825 & 46.704966597550765 & 50.1168617642713 & 49.685001919376916 & 38.206107919464145 \\
    & 4  & 55.54039002085877 & 58.022387334887334 & 60.5360833123991 & 60.76066313520357 & 41.82590392654731 \\
    & 8  & 61.41135687229437 & 66.19184027777779 & 67.0110367063492 & 66.9660330988456 & 47.78336998649498 \\
    & 16 & 63.81056472462723 & 72.49340879028378 & 73.50190145502646 & 73.33304398148148 & 56.986363462925965 \\
    \bottomrule
    \end{tabular}%
    } % Fine resizebox
    \label{tab:artgraph_results_precision_macro}
\end{table}

\begin{table}[H]
    \centering
    \caption{Risultati Medi di Recall Macro (\%) su ArtGraph. Ogni valore rappresenta la media dei punteggi di Recall Macro ottenuti su 10 diversi seed.}
    \sisetup{
      round-mode=places,
      round-precision=2,
      output-decimal-marker={,},
      table-format=2.2 
    }
    \resizebox{\textwidth}{!}{%
    \begin{tabular}{@{}lcSSSSS@{}}
    \toprule
    Backbone & Shots & {\rotatebox{90}{APE}} & {\rotatebox{90}{GDA\_CLIP}} & {\rotatebox{90}{TIMO}} & {\rotatebox{90}{TIMO\_S}} & {\rotatebox{90}{Tip\_Adapter}} \\
    \midrule
    \multirow{5}{*}{RN101}
    & 1  & 42.03125 & 41.25868055555555 & 45.47743055555556 & 45.46875 & 38.081597222222214 \\
    & 2  & 47.421875 & 46.31944444444444 & 46.979166666666664 & 48.68923611111111 & 40.77256944444444 \\
    & 4  & 52.41319444444444 & 55.3125 & 57.46527777777777 & 57.78645833333333 & 45.12152777777777 \\
    & 8  & 57.25694444444444 & 63.628472222222214 & 64.22743055555556 & 64.64409722222221 & 49.52256944444444 \\
    & 16 & 58.97569444444444 & 70.72916666666666 & 70.52083333333334 & 70.703125 & 54.07118055555556 \\
    \midrule
    \multirow{5}{*}{RN50}
    & 1  & 43.94097222222222 & 40.76388888888889 & 43.51562499999999 & 45.53819444444444 & 38.37673611111111 \\
    & 2  & 46.64930555555556 & 46.41493055555555 & 47.31770833333333 & 47.951388888888886 & 40.32986111111111 \\
    & 4  & 53.71527777777777 & 54.96527777777777 & 55.90277777777777 & 56.744791666666664 & 44.375 \\
    & 8  & 57.96875 & 63.19444444444444 & 64.93055555555556 & 64.74826388888889 & 47.99479166666667 \\
    & 16 & 61.953125 & 72.578125 & 72.94270833333334 & 73.03819444444444 & 52.369791666666664 \\
    \midrule
    \multirow{5}{*}{ViT-B/16}
    & 1  & 50.416666666666664 & 48.26388888888889 & 51.744791666666664 & 52.760416666666664 & 46.467013888888886 \\
    & 2  & 54.35763888888889 & 54.887152777777786 & 53.68923611111111 & 55.503472222222214 & 48.44618055555556 \\
    & 4  & 59.947916666666664 & 62.760416666666664 & 63.73263888888889 & 64.23611111111111 & 52.005208333333336 \\
    & 8  & 63.550347222222214 & 70.46006944444444 & 71.09375 & 71.06770833333333 & 57.734375 \\
    & 16 & 68.125 & 77.34375 & 77.75173611111111 & 77.56944444444443 & 61.91840277777777 \\
    \midrule
    \multirow{5}{*}{ViT-B/32}
    & 1  & 44.947916666666664 & 42.717013888888886 & 46.796875 & 47.38715277777777 & 40.95486111111111 \\
    & 2  & 47.717013888888886 & 46.5625 & 48.72395833333333 & 48.697916666666664 & 42.10069444444444 \\
    & 4  & 54.51388888888889 & 56.03298611111111 & 58.68923611111111 & 58.784722222222214 & 45.347222222222214 \\
    & 8  & 59.583333333333336 & 63.090277777777786 & 64.14930555555556 & 64.08854166666667 & 49.40104166666667 \\
    & 16 & 62.369791666666664 & 69.74826388888889 & 70.76388888888889 & 70.72048611111111 & 56.041666666666664 \\
    \bottomrule
    \end{tabular}%
    } % Fine resizebox
    \label{tab:artgraph_results_recall_macro}
\end{table}

\begin{table}[H]
    \centering
    \caption{Risultati Medi di F1-Score Macro (\%) su ArtGraph. Ogni valore rappresenta la media dei punteggi di F1-Score Macro ottenuti su 10 diversi seed.}
    \sisetup{
      round-mode=places,
      round-precision=2,
      output-decimal-marker={,},
      table-format=2.2 
    }
    \resizebox{\textwidth}{!}{%
    \begin{tabular}{@{}lcSSSSS@{}}
    \toprule
    Backbone & Shots & {\rotatebox{90}{APE}} & {\rotatebox{90}{GDA\_CLIP}} & {\rotatebox{90}{TIMO}} & {\rotatebox{90}{TIMO\_S}} & {\rotatebox{90}{Tip\_Adapter}} \\
    \midrule
    \multirow{5}{*}{RN101}
    & 1  & 38.22480832517401 & 35.92247954009474 & 41.82578425455264 & 41.38888866629744 & 32.543459168991355 \\
    & 2  & 44.15010495593747 & 42.21386966248248 & 44.38232820885685 & 45.60103994249215 & 35.061724463608776 \\
    & 4  & 49.124615475394094 & 52.796017736828766 & 55.09538485322154 & 55.41959194808459 & 39.507916497388386 \\
    & 8  & 53.94211276912515 & 61.86328598552863 & 62.6860992017242 & 63.04887936137936 & 43.94208926205231 \\
    & 16 & 55.59800859203433 & 69.58837574462575 & 69.40211524586525 & 69.57427641802641 & 49.02615712230566 \\
    \midrule
    \multirow{5}{*}{RN50}
    & 1  & 39.75858247986247 & 35.32592723652289 & 40.291962922018065 & 41.068853859594746 & 30.917586834729537 \\
    & 2  & 42.37257063040732 & 42.46836105001656 & 43.723073312257746 & 43.98020758892953 & 32.581819715473515 \\
    & 4  & 49.72670853363388 & 52.661093502086146 & 53.95745567427068 & 54.54884865270894 & 37.18713264738134 \\
    & 8  & 54.03185332505185 & 61.57093196155696 & 63.66344477281977 & 63.454451220076216 & 41.19705450508932 \\
    & 16 & 58.358082311207305 & 71.62958627021126 & 72.09026308658662 & 72.10209673904527 & 46.3837157863135 \\
    \midrule
    \multirow{5}{*}{ViT-B/16}
    & 1  & 46.86134621487735 & 43.44192866754061 & 48.77821810452905 & 49.376949866300784 & 40.265990582989424 \\
    & 2  & 51.217461035016186 & 51.642804820702466 & 50.99345978436052 & 52.442428503549834 & 42.90267840819312 \\
    & 4  & 57.061529981349054 & 60.658542964173776 & 61.60608026233026 & 62.13342964905464 & 46.927358946200854 \\
    & 8  & 60.74658443408443 & 69.19408542846043 & 69.86315767565767 & 69.7550609113109 & 53.409287271322874 \\
    & 16 & 65.6820766039516 & 76.49890445202945 & 76.9664797008547 & 76.60665607540606 & 58.83030334317099 \\
    \midrule
    \multirow{5}{*}{ViT-B/32}
    & 1  & 40.71803339992194 & 36.725527141703616 & 42.78675947232448 & 42.98310193346958 & 34.796268284948624 \\
    & 2  & 43.91984343271108 & 41.871811822687405 & 45.341659391412676 & 45.2947052177412 & 35.742660825823094 \\
    & 4  & 50.561446845817 & 52.93884228856656 & 55.79287928506678 & 55.84903695841196 & 39.00138294439765 \\
    & 8  & 56.005300833425835 & 61.37481430266357 & 62.30416747604247 & 62.29904297091797 & 43.660447463756284 \\
    & 16 & 58.52284058534059 & 67.89625652125652 & 69.09418532856033 & 68.97160709660709 & 51.17863151962416 \\
    \bottomrule
    \end{tabular}%
    } % Fine resizebox
    \label{tab:artgraph_results_f1_macro}
\end{table}

\subsection{Analisi dei Risultati}
L'analisi dei risultati medi, ottenuti mediando le performance su 10 diversi seed per ciascuna configurazione, permette di delineare le seguenti osservazioni chiave:

\begin{enumerate}
    \item \textbf{Impatto del Numero di Shots (Campioni di Addestramento):}
    \begin{itemize}[label=\textbullet, leftmargin=*]
        \item Si osserva un incremento consistente delle performance (Accuratezza, Precisione Macro, Recall Macro, F1-Score Macro) per tutti i backbone e modelli di adattamento all'aumentare del numero di ``shots''. Questo conferma l'efficacia di fornire più esempi specifici per il task target.
        \item I guadagni di performance sono generalmente più marcati nei passaggi iniziali (es. da 1 a 2 shots, o da 2 a 4 shots) rispetto agli incrementi con un numero maggiore di campioni (es. da 8 a 16 shots), suggerendo rendimenti decrescenti ma comunque positivi.
    \end{itemize}

    \item \textbf{Performance dei Backbone Pre-addestrati:}
    \begin{itemize}[label=\textbullet, leftmargin=*]
        \item Il backbone \texttt{ViT-B/16} emerge come il più performante in maniera predominante su tutte le metriche, raggiungendo i valori più elevati in quasi tutte le condizioni.
        \item \texttt{ViT-B/32} si posiziona generalmente come il secondo miglior backbone, superando i modelli ResNet ma rimanendo al di sotto di \texttt{ViT-B/16}.
        \item I backbone \texttt{RN50} e \texttt{RN101} mostrano prestazioni competitive tra loro, ma sono tendenzialmente superati dai modelli ViT.
    \end{itemize}

    \item \textbf{Efficacia dei Modelli di Adattamento Few-Shot (basata sull'Accuratezza):}
    \begin{itemize}[label=\textbullet, leftmargin=*]
        \item \texttt{Tip\_Adapter} registra costantemente le prestazioni di accuratezza più basse, fungendo da baseline ma venendo nettamente superato dagli altri metodi.
        \item \texttt{APE} mostra un miglioramento rispetto a \texttt{Tip\_Adapter}, ma si colloca in una posizione intermedia.
        \item I modelli \texttt{GDA\_CLIP}, \texttt{TIMO} e \texttt{TIMO\_S} si dimostrano i più efficaci in termini di accuratezza, contendendosi le prime posizioni.
            \begin{itemize}[label=\textendash, leftmargin=*]
                \item \textit{Con pochi shots (1-shot):} \texttt{TIMO\_S} e \texttt{TIMO} tendono a mostrare un leggero vantaggio o a essere i migliori (es. \texttt{TIMO\_S} con \texttt{ViT-B/16} raggiunge \SI{53.58}{\percent} di accuratezza).
                \item \textit{Con un numero maggiore di shots (specialmente 8 e 16):} Le performance di accuratezza di \texttt{GDA\_CLIP}, \texttt{TIMO} e \texttt{TIMO\_S} diventano molto ravvicinate. \texttt{TIMO} con \texttt{ViT-B/16} e 16 shots ottiene la performance di accuratezza più alta in assoluto (\SI{78.10}{\percent}), seguito da vicino da \texttt{TIMO\_S} e \texttt{GDA\_CLIP}.
            \end{itemize}
        \item La differenza di accuratezza tra \texttt{TIMO} e \texttt{TIMO\_S} è generalmente minima, con \texttt{TIMO\_S} che a volte presenta un piccolo margine con un numero esiguo di dati.
    \end{itemize}
    
    \item \textbf{Analisi delle Metriche di Precisione Macro, Recall Macro e F1-Score Macro:}
    \begin{itemize}[label=\textbullet, leftmargin=*]
        \item \textbf{Tendenze Generali Consistenti:} Le metriche di Precisione Macro, Recall Macro e F1-Score Macro seguono tendenze generali molto simili a quelle osservate per l'accuratezza. Si registra un miglioramento delle prestazioni con l'aumento del numero di shots. Il backbone \texttt{ViT-B/16} si conferma generalmente il più performante, e i metodi di adattamento \texttt{TIMO}, \texttt{TIMO\_S} e \texttt{GDA\_CLIP} rimangono i leader, mentre \texttt{Tip\_Adapter} è costantemente il meno performante e \texttt{APE} si posiziona a un livello intermedio.
        \item \textbf{Precisione vs. Recall:} Per molti dei metodi e backbone più performanti, specialmente con un numero elevato di shots, i valori di Precisione Macro tendono ad essere leggermente superiori a quelli di Recall Macro. Ad esempio, con \texttt{ViT-B/16} e 16 shots, \texttt{TIMO} ottiene una Precisione Macro del \SI{80.62}{\percent} e un Recall Macro del \SI{77.75}{\percent}. Questo suggerisce che i modelli sono, in media, leggermente più conservativi, preferendo evitare falsi positivi a costo di qualche falso negativo in più.
        \item \textbf{F1-Score Macro come Bilanciamento:} L'F1-Score Macro, che bilancia Precisione e Recall, vede anch'esso \texttt{TIMO}, \texttt{TIMO\_S} e \texttt{GDA\_CLIP} come i metodi migliori. Con \texttt{ViT-B/16} e 16 shots, \texttt{TIMO} raggiunge l'F1-Score Macro più elevato (\SI{76.97}{\percent}), seguito da \texttt{TIMO\_S} (\SI{76.61}{\percent}) e \texttt{GDA\_CLIP} (\SI{76.50}{\percent}). Questi valori sono tipicamente leggermente inferiori alle rispettive accuratezze quando precisione e recall non sono perfettamente bilanciate.
        \item \textbf{Performance Relative di TIMO e TIMO\_S:} Le differenze tra \texttt{TIMO} e \texttt{TIMO\_S} rimangono minime anche per queste metriche, con variazioni marginali a seconda della configurazione specifica, rispecchiando quanto osservato per l'accuratezza.
        \item \textbf{Gerarchia dei Metodi:} La gerarchia generale delle prestazioni dei metodi di adattamento (\texttt{Tip\_Adapter} $<$ \texttt{APE} $<$ \{\texttt{GDA\_CLIP}, \texttt{TIMO}, \texttt{TIMO\_S}\}) è sostanzialmente mantenuta attraverso tutte le metriche considerate.
    \end{itemize}

    \item \textbf{Combinazioni Ottimali e Performance Massime (Considerando Tutte le Metriche):}
    \begin{itemize}[label=\textbullet, leftmargin=*]
        \item Le performance più elevate su tutte le metriche si ottengono combinando il backbone \texttt{ViT-B/16} con 16 shots e i modelli di adattamento \texttt{TIMO}, \texttt{TIMO\_S} o \texttt{GDA\_CLIP}. Ad esempio, \texttt{TIMO} con \texttt{ViT-B/16} a 16 shots raggiunge Accuratezza del \SI{78.10}{\percent}, Precisione Macro del \SI{80.62}{\percent}, Recall Macro del \SI{77.75}{\percent} e F1-Score Macro del \SI{76.97}{\percent}.
        \item Anche con 8 shots, la combinazione di \texttt{ViT-B/16} con questi tre modelli di adattamento permette di ottenere risultati molto competitivi su tutte le metriche (es. F1-Score Macro intorno al \SI{69}{\percent}-\SI{70}{\percent}).
    \end{itemize}

    \item \textbf{Considerazioni Generali:}
    \begin{itemize}[label=\textbullet, leftmargin=*]
        \item La mediazione dei risultati su 10 diversi seed fornisce una stima più robusta e affidabile delle prestazioni, riducendo l'impatto della variabilità stocastica.
        \item Emerge chiaramente la sinergia tra un backbone potente (come \texttt{ViT-B/16}) e un metodo di adattamento few-shot efficace per massimizzare le performance su un ampio spettro di metriche.
        \item Nonostante i miglioramenti, il dataset \texttt{artgraph} si conferma sfidante, con le migliori performance che indicano margini per ulteriori progressi, specialmente nel bilanciamento tra precisione e recall.
    \end{itemize}
\end{enumerate}
\section{Proposta: StyTIMO}
StyTIMO è un'estensione proposta di TIMO che integra informazioni stilistiche nelle rappresentazioni visive, particolarmente rilevante per la classificazione di opere d'arte. Mentre TIMO si concentra sulla guida reciproca tra rappresentazioni testuali e visive semantiche, StyTIMO aggiunge un terzo elemento: la rappresentazione esplicita dello stile artistico.

\subsection{Motivazione}
La classificazione di opere d'arte per artista è una sfida particolare poiché richiede di cogliere non solo il contenuto semantico dell'opera (cosa rappresenta), ma anche le caratteristiche stilistiche distintive dell'artista (come è rappresentato). I modelli pre-addestrati come CLIP tendono a focalizzarsi principalmente sulle caratteristiche semantiche, perdendo potenzialmente le sottili differenze stilistiche che caratterizzano gli artisti.

\subsection{Architettura di StyTIMO}
StyTIMO estende TIMO integrando caratteristiche stilistiche esplicite attraverso un modulo chiamato StyleAdapter, che cattura e incorpora le informazioni dello stile artistico nelle rappresentazioni visive. L'architettura si articola in quattro componenti principali:

\begin{enumerate}
    \item \textbf{Estrazione di Feature Stilistiche}: Utilizza le matrici di Gram calcolate su feature map intermedie della rete CLIP per catturare le caratteristiche stilistiche delle immagini.
    
    \item \textbf{Proiezione delle Feature Stilistiche}: Riduce la dimensionalità delle matrici di Gram vettorizzate attraverso layer lineari addestrabili.
    
    \item \textbf{Fusion Adapter}: Integra le feature semantiche standard di CLIP con le feature stilistiche proiettate attraverso un adattatore con architettura bottleneck.
    
    \item \textbf{Pipeline TIMO Potenziata}: Utilizza le rappresentazioni visive arricchite di stile all'interno del framework di guida reciproca di TIMO.
\end{enumerate}

\subsection{Estrazione delle Feature Stilistiche con Matrici di Gram}
Le matrici di Gram sono ampiamente utilizzate nella letteratura sul trasferimento stilistico per catturare le caratteristiche statistiche dello stile di un'immagine:

\begin{enumerate}
    \item \textbf{Selezione di Layer}: Le feature map intermedie (tipicamente dai layer \texttt{layer2} e \texttt{layer3} per ResNet-50) vengono estratte tramite hook durante il forward pass.
    
    \item \textbf{Calcolo delle Matrici di Gram}: Per una feature map di forma $[B,C,H,W]$, la matrice di Gram è calcolata come:
    \[G = \frac{FF^T}{H \cdot W}\]
    dove $F$ è la feature map ridimensionata a $[C, H \cdot W]$.
    
    \item \textbf{Vettorizzazione Efficiente}: Dato che la matrice di Gram è simmetrica, viene vettorizzata estraendo solo la triangolare superiore, riducendo la dimensionalità da $C^2$ a $C(C+1)/2$.
\end{enumerate}

\subsection{Fusion Adapter con Architettura Bottleneck}
Il Fusion Adapter integra le informazioni semantiche con quelle stilistiche:

\begin{enumerate}
    \item \textbf{Concatenazione}: Le feature semantiche di CLIP e le feature stilistiche proiettate vengono concatenate.
    
    \item \textbf{Proiezione Bottleneck}: Il vettore concatenato è processato attraverso un'architettura bottleneck:
    \[input \rightarrow \text{down-projection} \rightarrow \text{layer norm} \rightarrow \text{ReLU} \rightarrow \text{dropout} \rightarrow \text{up-projection} \rightarrow output\]
    
    \item \textbf{Connessione Residuale Scalata}: Una connessione residuale modulata da un fattore $\alpha$ garantisce che le feature semantiche originali non vengano perse durante la fusione.
\end{enumerate}

\subsection{Formula di StyTIMO}
L'algoritmo di inferenza di StyTIMO può essere descritto come:

\begin{enumerate}
    \item \textbf{Estrazione Feature}: Per un'immagine query $q$:
    \begin{itemize}
        \item Estrai feature semantiche: $f_{\text{sem}}(q)$
        \item Estrai feature stilistiche dalle matrici di Gram: $G_l(q)$ per ogni layer $l$
        \item Vettorizza e proietta le matrici di Gram: $f_{\text{sty}}(q) = \text{concat}(P_l(G_l(q)))$
        \item Fondi le feature: $f_{\text{fused}}(q) = \text{FusionAdapter}(\text{concat}(f_{\text{sem}}(q), f_{\text{sty}}(q)))$
    \end{itemize}
    
    \item \textbf{Applicazione TIMO}: Utilizza $f_{\text{fused}}(q)$ al posto di $\fv$ standard nelle componenti TGI e IGT di TIMO.
\end{enumerate}

\subsection{Conclusione}

StyTIMO mira a superare TIMO nelle prestazioni di classificazione di opere d'arte, combinando efficacemente le capacità di guida mutua testo-immagine di TIMO con una rappresentazione esplicita dello stile artistico attraverso le matrici di Gram.

\newpage
\appendix
\section{Pseudo-codice di TIMO}
L'algoritmo 1 del paper originale fornisce uno pseudo-codice in stile PyTorch per l'algoritmo principale di TIMO.

\begin{algorithm}[H]
\caption{Algoritmo principale di TIMO}
\label{alg:timo}
\begin{algorithmic}[1]
\State \textbf{Input:} Feature immagini di supporto $\Fv \in \R^{N \times K \times D}$, feature testuali $\Ft \in \R^{N \times P \times D}$, feature immagine di query $\fv \in \R^D$, `ImageClassifierBuilder`, `TextClassifierBuilder`.
\State \textbf{Parametri:} Peso del branch multimodale $\alpha$, controllore del grado di fusione $\gamma_{\text{IGT}}$ (chiamato $\gamma$ nel paper per IGT, qui $\gamma_{\text{IGT}}$ per chiarezza), controllore grado di guida testuale $\beta_{\text{TGI}}$ (chiamato $\beta$ nel paper per TGI).
\Statex
\State \textit{// Modellazione}
\State $\mathbf{F}_{\text{vp}} \gets \Fv.\text{mean}(1).\text{unsqueeze}(1)$ \Comment{Prototipi visivi per classe}
\State $\mathbf{F}_{\text{vp}} \gets \mathbf{F}_{\text{vp}} / \mathbf{F}_{\text{vp}}.\text{norm}(-1, \text{keepdim=True})$ \Comment{Normalizzazione L2}
\State $\mathbf{W}_{\text{sim}} \gets \text{CosineSimilarity}(\Ft, \mathbf{F}_{\text{vp}}, -1)$ \Comment{Pesi di similarità per TGI e IGT}
\Statex
\State \textit{// Costruzione Classificatore TGI}
\State $\mathbf{S}_{\text{TGI}} \gets \text{SelezionaPesi}(\mathbf{W}_{\text{sim}}, \beta_{\text{TGI}})$ \Comment{Selezione/pesatura basata su $\beta_{\text{TGI}}$}
\State $\Ftgi \gets \text{concat}([\Fv, \Ft \odot \mathbf{S}_{\text{TGI}}], 1)$ \Comment{⊙ è prodotto element-wise o selezione}
\State $\text{Classifier}_{\text{TGI}} \gets \text{ImageClassifierBuilder}(\Ftgi)$
\Statex
\State \textit{// Costruzione Classificatore IGT}
\State $\mathbf{R}_{\text{IGT}} \gets \text{softmax}(\gamma_{\text{IGT}} \cdot \mathbf{W}_{\text{sim}} / \mathbf{W}_{\text{sim}}.\text{norm}(-1, \text{keepdim=True}), -1)$
\State $\Figt \gets \text{einsum("np, npd -> nd", } \mathbf{R}_{\text{IGT}}, \Ft \text{)}$ \Comment{Feature testuali rettificate}
\State $\text{Classifier}_{\text{IGT}} \gets \text{TextClassifierBuilder}(\Figt)$
\Statex
\State \textit{// Inferenza}
\State $\text{logits}_{\text{TGI\_out}} \gets \text{Classifier}_{\text{TGI}}(\fv)$
\State $\text{logits}_{\text{IGT\_out}} \gets \text{Classifier}_{\text{IGT}}(\fv)$
\State $\text{logits} \gets \text{logits}_{\text{TGI\_out}} + \alpha \cdot \text{logits}_{\text{IGT\_out}}$
\State \Return $\text{logits}$
\end{algorithmic}
\end{algorithm}
\textit{Nota: Lo pseudo-codice fornito nel paper è una versione semplificata. Ad esempio, $\mathbf{W}_{\text{sim}}$ qui rappresenta i pesi $s^i$ per TGI e i pesi $r^i$ (prima della softmax) per IGT. $\beta_{\text{TGI}}$ e $\gamma_{\text{IGT}}$ controllano il modo in cui questi pesi vengono usati.}

\end{document}